% Generated draft: requires mathematical and source review.
\documentclass[11pt]{article}
\usepackage[margin=25mm]{geometry}
\usepackage{fontspec}
\setmainfont{Latin Modern Roman}
\usepackage{amsmath,amssymb,mathtools,unicode-math}
\setmathfont{Latin Modern Math}
\usepackage{enumitem,needspace,xurl,hyperref}
\setlength{\parindent}{0pt}
\setlength{\parskip}{5pt}
\setlength{\emergencystretch}{3em}
\newcommand{\N}{\mathbb{N}}
\newcommand{\Z}{\mathbb{Z}}
\newcommand{\Q}{\mathbb{Q}}
\newcommand{\R}{\mathbb{R}}
\newcommand{\C}{\mathbb{C}}
\newcommand{\F}{\mathbb{F}}
\newcommand{\A}{\mathbb{A}}
\newcommand{\PP}{\mathbb{P}}
\newcommand{\E}{\mathbb{E}}
\newcommand{\eps}{\varepsilon}
\newcommand{\dd}{\,\mathrm{d}}
\begin{document}
\section*{UnsolvedMath ID: 30002385}
\textbf{Source problem code:} catalogue entry
\subsection{Definitions and mathematical statement}
The supplied problem concerns the isospectral torus \(T_E\) associated with a spectral set \(E\), but does not specify the admissible sets \(E\), define \(T_E\), or identify the underlying operator realization. Consequently, a fully quantified statement requires additional hypotheses. The following makes the usual coefficient interpretation explicit conditionally, without asserting that these missing conventions have been established by the source record.

Suppose a self-adjoint Jacobi-operator setting has been fixed, including its index set \(I\), domain, and any boundary condition. Write \(J=J(a,b)\), where \(a_n>0\) are the off-diagonal coefficients and \(b_n\in\R\) are the diagonal coefficients. For \(J_0=J(a^{(0)},b^{(0)})\in T_E\), interpret a coefficient \(\ell^2\) perturbation by
\[
 \sum_{n\in I}\left(\lvert a_n-a_n^{(0)}\rvert^2+\lvert b_n-b_n^{(0)}\rvert^2\right)<\infty,
\]
with the off-diagonal sum understood to range over the bonds present in the chosen realization. Define the corresponding conditional class
\[
 \mathcal P_2(T_E)=\left\{J(a,b):\ \exists J(a^{(0)},b^{(0)})\in T_E\ \text{with the displayed sum finite}\right\}.
\]
For a specified unit vector \(u\), its scalar spectral measure is the Borel probability measure
\[
 \mu_{J,u}(B)=\langle u,\mathbf 1_B(J)u\rangle
 \qquad(B\subseteq\R\ \text{Borel}),
\]
where \(\mathbf 1_B(J)\) is the spectral projection of \(J\). Write its Lebesgue decomposition as
\[
 \dd\mu_{J,u}(x)=w_{J,u}(x)\,\dd x+\dd\mu_{J,u,\mathrm s}(x).
\]

\textbf{Requested characterization.} Once the missing operator and spectral-measure conventions are fixed, give necessary and sufficient conditions for a measure to be the spectral measure of some \(J\in\mathcal P_2(T_E)\). The quantifier ranges over all base points \(J_0\in T_E\) and all admissible coefficient perturbations satisfying the displayed summability condition. If the intended setting requires matrix-valued spectral measures, the target must be adjusted accordingly.

\textbf{Requested preservation question.} Determine whether every such perturbation preserves the absolutely continuous spectrum of its base operator. One possible precise reading is
\[
 \forall J_0\in T_E\ \forall J\text{ satisfying the displayed condition relative to }J_0,
 \qquad \sigma_{\mathrm{ac}}(J)=\sigma_{\mathrm{ac}}(J_0),
\]
where \(\sigma_{\mathrm{ac}}(J)\) denotes the spectrum of the restriction of \(J\) to its absolutely continuous spectral subspace. The source does not specify whether this equality, preservation of spectral multiplicity, or equivalence of absolutely continuous spectral measure classes is intended. These formulations must not be identified without further hypotheses. In particular, the supplied record does not justify imposing a finite-gap hypothesis on \(E\), choosing a half-line or whole-line realization, or selecting a distinguished spectral vector.
\subsection{Short English statement}
Characterize the spectral measures of every \(\ell^2\) perturbation of every point of the isospectral torus \(T_E\), and determine whether these perturbations preserve absolutely continuous spectrum. The source leaves the hypotheses on \(E\), the operator realization, the spectral-measure convention, and the exact meaning of preservation unspecified. The displayed formulation is therefore conditional and requires human mathematical review.
\subsection{Sources}
\begin{itemize}
\item Supplied local UnsolvedMath registry record, permanent ID \(30002385\), title \emph{Spectral Measures of Perturbed Isospectral Tori}, source problem code \texttt{catalogue entry}, in the Oberwolfach Reports Open Problems collection. No individual proposer is identified; the supplied proposed year is 2013.
\item Supplied source citation: \emph{Mini-Workshop: Direct and Inverse Spectral Theory of Almost Periodic Operators} (2014), Oberwolfach Reports, pp.~2083--2117, DOI \texttt{10.4171/owr/2013/36}. The literature-source label is \emph{Oberwolfach Report 36/2013}. \url{https://doi.org/10.4171/owr/2013/36}. The discrepancy between the supplied 2014 citation year and the 2013 report designation is unresolved.
\item Supplied literature source: \emph{Finite Gap Jacobi Matrices III}. \url{https://arxiv.org/abs/1108.0183}. The record supplies no authors, theorem numbers, or passages establishing its precise relationship to the requested characterization.
\item Literature assessment dated exactly \texttt{2026-08-21}: the supplied curation labels the problem open and reports that, despite known finite-gap Szeg\H{o} and perturbation theorems, a general characterization for all \(\ell^2\) perturbations, including preservation of absolutely continuous spectrum, remains open. This is supplied dated curation, not an independently verified literature assessment. The background expressly requests independent verification of time-sensitive status and cited claims.
\item The registry labels the statement \texttt{corrected\_verified} and says that notation, typography, and wording were normalized without changing the mathematical content. These are supplied editorial assertions, not independent verification of the formulation or its mathematical status.
\end{itemize}
\end{document}
